\documentclass[10pt,twocolumn]{article}

\usepackage[margin=0.75in]{geometry}
\usepackage{amsmath}
\usepackage{booktabs}
\usepackage{xcolor}
\usepackage{graphicx}
\usepackage[hidelinks]{hyperref}
\usepackage{url}

\newcommand{\sys}{AgentTx}
\newcommand{\todo}[1]{\textcolor{red}{[\textbf{TODO:} #1]}}

\title{\vspace{-1.5em}\textbf{Beyond Compensation: Kernel-Level Deferral of\\
Irreversible Effects for Autonomous AI Agents}\\[0.3em]
\large The \sys{} Project Proposal}

\author{
  \textit{Author One} \and \textit{Author Two} \and
  \textit{Author Three} \and \textit{Author Four}\\[0.3em]
  \small Department of Computer Science and Engineering\\
  \small \todo{Institution} \quad \todo{Course code, B.Tech.\ Year III}
}
\date{\today}

\begin{document}
\maketitle

\begin{abstract}
Filesystem damage caused by autonomous coding agents is a solved problem.
Production agents checkpoint every file-modifying action and restore on demand,
and execute inside Landlock, seccomp and Seatbelt sandboxes by default. What
remains unsolved --- and is publicly conceded as \emph{unsolvable} --- is the
outbound effect: the push, the webhook, the email, the payment. Once a third
party has observed an action, the prevailing view holds, no mechanism can take
it back.

That view is correct but stated at the wrong boundary. Irreversibility begins at
\emph{observation}, not at the system call. Every deployed remedy ---
compensating transactions, sagas, idempotency keys --- operates downstream of
emission and is therefore best-effort by construction. Between the agent's
\texttt{sendmsg} and the packet leaving the machine lies a region nobody
occupies: the kernel.

We propose \sys{}, which relocates the commit point into that region. Writes
divert into a copy-on-write layer and outbound effects are captured in a kernel
write-ahead log rather than emitted; commit applies both, abort discards both.
For fire-and-forget effects this converts best-effort compensation into a hard
guarantee. We do not reverse observation --- we prevent it.

The mechanism, however, is only as meaningful as the policy governing it, and
the governing question is \emph{what constitutes a transaction for an agent}.
Too short and the guarantee window is negligible; too long and the system holds
an unbounded write-ahead log across an unbounded number of agent decisions,
a regime the transactional-memory literature was never designed for. We argue
for \emph{verification-delimited} transactions: a transaction spans the interval
over which no evidence about the agent's work is available, and commits the
moment evidence arrives.

\sys{} contributes (i) kernel-level deferral of fire-and-forget external
effects; (ii) verification-delimited transaction boundaries, with the
guarantee-window against staleness trade-off measured rather than asserted;
(iii) an \texttt{int8}-quantised classifier evaluated \emph{inside} an eBPF LSM
hook; (iv) a four-class effect taxonomy grounded in the literature on
irrevocability, with the first measured distribution of real agent traces across
it; and (v) the first head-to-head evaluation of agent rollback mechanisms under
a single workload and metric set.
\end{abstract}

% =====================================================================
\section{Introduction}

A prompt injection is planted in the documentation of a widely used dependency.
An agent reads it in the course of an ordinary task, and --- acting with the
user's full authority --- reads the user's SSH private key and POSTs it to an
external host. The POST is fire-and-forget: the agent does not require the
response in order to continue.

Under every deployed protection this either succeeds, because the destination
allowlist happened to permit the host, or because the approval prompt displayed
a shell command whose system-call trace the user could not have predicted; or it
is blocked by a static rule that somebody wrote in advance having anticipated
this specific exfiltration path. Neither is a guarantee.

From the kernel's perspective the agent is an ordinary process running under the
user's UID. From the user's perspective it is a process whose behaviour is
partly determined by text it read moments earlier --- a confused deputy that is
non-deterministic and ingests attacker-influenced input as a matter of routine
operation.

\subsection{What is already solved}

We concede a great deal at the outset, because the value of this proposal
depends on being precise about what remains.

\paragraph{Sandboxing is commodity.} Codex CLI executes commands inside
bubblewrap with unprivileged user namespaces, Landlock LSM filesystem filtering
and seccomp-BPF syscall restriction on Linux, Seatbelt on macOS, and restricted
tokens with synthetic SIDs on Windows~\cite{codexsandbox}. Anthropic has
open-sourced the isolation layer behind Claude Code's sandbox
tooling~\cite{sandboxescapes}.

\paragraph{Filesystem rollback is commodity.} Claude Code records a checkpoint
of affected files before every filesystem-modifying action and exposes a rewind
operation restoring any prior point in a session~\cite{ccrewind}. The research
literature has converged on the same capability~\cite{ftsandbox2025,
deltabox2026}.

\paragraph{Preventive egress control ships today.} Codex applies default-deny
egress with destination allowlists, performing DNS and IP classification before
permitting a hostname~\cite{codexsandbox}. ChatGPT Work exposes plan approval,
configurable check-ins and per-action approval
gates~\cite{chatgptwork,agentapprovals}.

\subsection{What is not solved}

Practitioners now describe agent reversibility in three tiers: the filesystem,
cheaply reversible through checkpoints and worktrees; the database, reversible
only through explicitly coded compensating transactions; and external side
effects --- a sent email, a captured payment, a delivered webhook --- described
as \emph{actions no mechanism can take back once a third party has seen
them}~\cite{rollbackpatterns2026}.

That third tier is a stated impossibility claim, and on its own terms it is
true. Observation transfers information, and no operation removes information
from a party one does not control. A refund is not an un-charge; it is a second
effect that happens to net to zero. Two-phase commit works only because
participants agree in advance to prepare; a third party who never joined the
protocol offers no lever. Speculator~\cite{nightingale2005speculator} and
Xsyncfs~\cite{nightingale2006xsyncfs} met this wall two decades ago with
considerable kernel machinery and drew the same line: they buffered output until
it was safe, and never claimed to un-output.

\subsection{The observation}

The claim conflates two distinct events.

\begin{center}
\small
\begin{tabular}{@{}cl@{}}
\multicolumn{2}{@{}l}{\texttt{agent calls sendmsg()}}\\
$\downarrow$ & \\
\textbf{1.\ Emission} & the packet leaves the machine\\
$\downarrow$ & \textit{$\leftarrow$ true irreversibility boundary}\\
\textbf{2.\ Observation} & a third party reads and acts\\
\end{tabular}
\end{center}

Compensating transactions, sagas and idempotency keys all begin downstream of
emission, and are best-effort by construction because they operate where the
impossibility genuinely binds. The interval between the system call and emission
--- wholly inside the kernel --- is unoccupied, and is the only place a
\emph{guarantee} rather than a best effort is available.

We note that this differs from idempotency keys in kind, not merely in degree:
an idempotency key requires the \emph{receiver} to implement it, whereas
deferral requires nothing of the receiver at all.

\subsection{Scope: what the mechanism does not reach}
\label{sec:scope}

Not every agent incident is within reach, and we state the boundary early
because a motivating example that the mechanism does not address is worse than
no example. In 2026 an agent deleted a production database and its backups in
nine seconds~\cite{cursorincident}. \sys{}'s \emph{classifier} would not have
caught it: each individual operation was routine, and per-call classification is
by construction a local decision. \sys{}'s \emph{deferral} would not have held
it either, since a remote database mutation is a request--response exchange in
which the agent blocks on the reply (\S\ref{sec:reqresp}).

What \sys{} offers in that scenario is the \emph{transaction}: a destructive
sequence staged under a verification-delimited boundary whose commit gate is
``the test suite still passes'' never commits, regardless of how routine each
step appeared. Mechanism and motivating example must be matched, and we
distinguish throughout which of our contributions addresses which failure.

\subsection{Contributions}

\begin{enumerate}\itemsep2pt
\item \textbf{Kernel-level deferral of fire-and-forget external effects.}
      Outbound operations issued inside a transaction are captured in a kernel
      write-ahead log and released only on commit. For this class, abort
      provides a hard guarantee where the state of the art provides best-effort
      compensation.
\item \textbf{Verification-delimited transaction boundaries.} We argue that the
      boundary question is the governing design decision, propose an answer
      derived from when evidence becomes available, and measure the resulting
      trade-off between guarantee window, snapshot staleness and log growth
      (\S\ref{sec:policy}).
\item \textbf{In-kernel effect classification.} An \texttt{int8}-quantised
      classifier executes with integer arithmetic inside an eBPF LSM hook,
      removing the per-call context switch that makes \texttt{seccomp}
      user-notify mediation too expensive to leave enabled.
\item \textbf{A four-class effect taxonomy}, grounded in the irrevocable
      transaction literature, with the first measured distribution of real agent
      system-call traces across it.
\item \textbf{A comparative evaluation} of agent rollback mechanisms under one
      workload and one metric set.
\end{enumerate}

\subsection{What we do not claim}
\label{sec:nonclaims}

We do not claim the transaction abstraction, which is
TxOS~\cite{porter2009txos}; output buffering, which is
Xsyncfs~\cite{nightingale2006xsyncfs}; fixed-point inference under the eBPF
verifier, which is established
practice~\cite{ebpffixedpoint2024,kernelpacketml2024,akgun2023kml}; filesystem
rollback for agents, which ships in production and appears in concurrent
work~\cite{ccrewind,ftsandbox2025,deltabox2026}; or the ability to reverse
observed effects, which is impossible.

We further decline to claim automatic inference of compensation windows
(\S\ref{sec:compensable}), which we believe the kernel cannot obtain and which
we name as an open problem rather than a contribution.

% =====================================================================
\section{Background and Related Work}
\label{sec:related}

\subsection{Operating system transactions}

TxOS~\cite{porter2009txos} added \texttt{sys\_xbegin}, \texttt{sys\_xend} and
\texttt{sys\_xabort} to Linux 2.6.22, providing ACID semantics for OS resources
using lazy version management over shadow copies of kernel objects. It cost
3{,}300 lines for transaction management and 5{,}300 for object management, plus
roughly 14{,}000 lines of changes to existing kernel code; a transactional
OpenSSH installation incurred 10\% overhead. It was never merged.

Its coverage is the relevant fact for us. TxOS provided transactional semantics
for \textbf{150 of 303 Linux system calls}, with 153 explicitly unsupported
(\cite[Table~1]{porter2009txos}). \texttt{socket} is among the unsupported
communication calls, alongside \texttt{ipc}, \texttt{mq\_open} and
\texttt{mq\_unlink}; the authors note that ``extending transactions to include
networking (a real challenge) would increase the count of supported calls by
5.'' On encountering an unsupported call, ``the system logs a warning or aborts
the transaction, depending on the flags passed to \texttt{sys\_xbegin()}''
(\cite[\S3.2]{porter2009txos}). There is no irrevocable-transaction mechanism:
un-undoable work is handled by exclusion, then by warning or abort --- never by
committing early under a no-abort guarantee. Nor does TxOS restore application
state on abort, only kernel state (\cite[\S3.1.4]{porter2009txos}).

Most directly relevant, TxOS anticipated our central limit. Its \S3.1.5
observes that ``code that communicates outside of a transaction and requires a
response cannot be encapsulated into a single system transaction,'' that ``the
system might buffer the message until commit,'' and that a transaction awaiting
a reply on the same channel will deadlock. Its resolution: ``\emph{the
programmer is responsible} for avoiding this send/reply idiom within a
transaction.''

TxFS~\cite{hu2018txfs,hu2019txfs} subsequently showed that ACID filesystem
transactions are obtainable in 5{,}200 lines by reusing the ext4 journal.

\emph{Our position.} The abstraction is inherited, and so --- we discovered on
reading it closely --- is the observation that request--response communication
cannot be transactional and that outbound messages may be buffered until commit.
What does not transfer is TxOS's resolution. \textbf{In the agent setting there
is no programmer} to be held responsible: the sequence of system calls is
generated at run time, non-deterministically, by a process under adversarial
influence. An escape hatch that reads as reasonable engineering advice in 2009
is precisely the assumption that fails in 2026. Our contribution is therefore
not the idea of buffering but its enforcement without a cooperating programmer:
classification of which effects may be buffered, at the network layer TxOS
excluded, with commit authority removed from the party issuing the calls
(\S\ref{sec:authority}).

\subsection{Irrevocability: the correct frame}

The problem of I/O inside a transaction is long-standing and named. The software
transactional memory literature calls a transaction performing un-undoable work
\emph{irrevocable} or \emph{inevitable}, and resolves it not by reversing the
I/O but by guaranteeing in advance that the transaction will never abort,
typically by serialising it so that it wins all
conflicts~\cite{welc2008irrevocable,spear2008inevitability}.

We adopt this frame, with one caveat we must confront rather than inherit: that
literature assumes short-lived transactions bounded by a few thousand
instructions. Agent transactions are bounded by human-scale activities lasting
seconds to minutes, and the assumption does not transfer.
\S\ref{sec:boundary} treats this as a first-class problem.

\subsection{Speculation and external output}

Speculator~\cite{nightingale2005speculator} provided kernel support for
speculative execution, checkpointing processes and propagating causal
dependencies through IPC. Xsyncfs~\cite{nightingale2006xsyncfs} developed
\emph{external synchrony}: buffer externally observable output until the state
it depends on is durable.

\emph{Our position.} Our effect log is external synchrony with a different
trigger --- we delay until the enclosing transaction commits rather than until
state is durable --- applied to network sends rather than file output, in a
setting where the emitting process is untrusted rather than merely impatient.
That last difference is not cosmetic: it forces a threat model
(\S\ref{sec:threat}) that neither prior system required.

\subsection{Agent sandboxing, rollback and state}

Concurrent work presents a transactional sandbox for coding agents combining
policy interception with filesystem snapshots, reporting complete interception
of high-risk commands and complete rollback at roughly 14.5\%
overhead~\cite{ftsandbox2025}. DeltaBox~\cite{deltabox2026} achieves millisecond
checkpoint and rollback through a change-based state abstraction.
Cordon~\cite{cordon2026} places transactions at the tool-call boundary.
AgentSight~\cite{agentsight2025} demonstrates eBPF observability at under 3\%
overhead by correlating intercepted model traffic with kernel events --- a
technique we adopt for payload inspection (\S\ref{sec:reqresp}).
ACRFence~\cite{acrfence2026} studies attacks against agent checkpoint-restore
itself. ChronoMem~\cite{chronomem2026} and MemTX~\cite{memtx2026} apply
versioning and transactional commit to agent \emph{memory}, which we draw on in
\S\ref{sec:coherence}. ActPlane~\cite{actplane} and CRAB~\cite{crab} are
reported to perform in-kernel eBPF policy enforcement and effect classification
respectively.
\todo{Read \cite{actplane} and \cite{crab} in full before the design freeze. If
either already defers external effects in-kernel, Contribution~1 collapses and
the project repositions onto Contributions~2 and~5.}

Orthogonally, CaMeL~\cite{debenedetti2025camel} defeats prompt injection by
construction. \sys{} is complementary: CaMeL constrains what an agent decides;
\sys{} constrains what its system calls can make real.

\subsection{Machine learning inside the kernel}

KML~\cite{akgun2023kml,akgun2021kmlarxiv} demonstrated in-kernel neural networks
and decision trees in under 4\,KB of dynamic memory at below 0.2\% CPU overhead.
Within eBPF the verifier forbids floating point; the established response is to
train in user space, quantise to integer parameters, ship weights through BPF
maps and execute layers via tail calls, with in-kernel inference roughly 84\%
faster than the user-space
equivalent~\cite{ebpffixedpoint2024,kernelpacketml2024,ebpfml2024}.

% =====================================================================
\section{Mechanism}
\label{sec:design}

\subsection{Interface}

\begin{center}
\texttt{tx\_begin(flags)}\quad\texttt{tx\_commit()}\quad\texttt{tx\_abort()}
\end{center}

\noindent During development these are \texttt{ioctl}s on a character device;
the final implementation promotes them to system calls. Critically,
\texttt{tx\_commit} is \emph{not} available to the transacting process
(\S\ref{sec:authority}).

\subsection{Isolation model}

On \texttt{tx\_begin} the transaction manager mounts an overlay in a private
mount namespace belonging to the transacting process: the working directory as
read-only lower layer, a per-transaction staging directory as upper, merged at
the original path and visible only inside that namespace.

The agent observes an unchanged path and reads its own writes. Every other
process --- the user's editor, shell and version control --- observes the
unmodified directory, yielding snapshot isolation and ensuring an abort cannot
invalidate an open editor buffer. Storage cost is proportional to the set of
files written, not the size of the working directory; where the underlying
filesystem supports cloning, copy-up may be satisfied by a reflink. We measure
both cases.

\subsection{The four-class taxonomy}

eBPF LSM programs attach to \texttt{file\_open}, \texttt{inode\_unlink},
\texttt{inode\_rename}, \texttt{socket\_connect}, \texttt{socket\_sendmsg} and
\texttt{bprm\_check\_security}. Each intercepted operation is classified:

\begin{center}
\small
\begin{tabular}{@{}lp{4.3cm}@{}}
\toprule
\textbf{Class} & \textbf{Handling}\\
\midrule
Reversible   & proceed; the CoW layer captures it\\
Deferrable   & record in the effect WAL, report success, do not emit\\
Compensable  & emit, record a cancellation window and commit deadline\\
Irrevocable  & the transaction becomes non-abortable and terminates; escalate
               to a human approval gate, deny on timeout\\
\bottomrule
\end{tabular}
\end{center}

Commit flushes the WAL in issue order and merges the upper layer down; abort
discards both. A write to a terminal has already been observed by a human and is
irrevocable by construction.

\subsection{The compensable class: a registry, not an inference}
\label{sec:compensable}

Some destinations do cooperate: a payment may be voided before capture, a push
force-reverted before anyone pulls, a resource deleted after creation. The
kernel, however, has no application semantics with which to discover this. We
considered and rejected inferring cancellation windows from traffic; nothing
observable at the LSM hook distinguishes an endpoint that honours a void from
one that does not.

We therefore specify the compensable class as a \emph{declarative registry},
distributed and updated like a certificate bundle: a mapping from host and path
pattern to a cancellation endpoint and a window duration. The kernel consults
it; it does not learn it. A transaction that emits a compensable effect acquires
a commit deadline equal to the shortest window it holds, and is aborted --- with
compensation attempted --- if the deadline passes.

This is deliberately modest. Automatic inference of compensation semantics is,
we believe, not solvable at the kernel boundary, and we name it as an open
problem rather than claim it (\S\ref{sec:nonclaims}).

\subsection{The request--response limit}
\label{sec:reqresp}

Deferral is possible only where the agent does not block on a reply. A
fire-and-forget effect --- a push, a webhook, a notification, a payment
submission --- may be held; a request--response exchange such as an API query or
package download may not, and a transaction that waits on a buffered channel
deadlocks outright.

This limit is not ours to discover: TxOS states it precisely in its \S3.1.5 and
resolves it by assigning the programmer responsibility for avoiding the
send/reply idiom inside a transaction. We inherit the limit and reject the
resolution, because no programmer is present to be held responsible. The
consequence is that what TxOS could leave as a documented caveat, we must detect
and enforce at run time: an effect that would block must be recognised as such
before it is buffered, or the agent deadlocks. Classification therefore does
work in \sys{} that it did not have to do in TxOS.

The magnitude of Contribution~1 is bounded by the fraction of agent outbound
effects that are fire-and-forget, a quantity nobody has measured. We treat this
as a first-class result and a project gate (\S\ref{sec:eval}).

A second limit is visibility: at the syscall boundary a TLS \texttt{sendmsg}
carries opaque bytes, so mutating and non-mutating operations to the same host
are indistinguishable. We adopt AgentSight's approach~\cite{agentsight2025},
attaching uprobes to the TLS library write path to recover plaintext. This is
precisely where \sys{} improves on destination allowlists, which are static and
per-host and cannot separate a read from a mutation against the same endpoint.

\subsection{In-kernel classification}

A small model is trained in user space on labelled agent traces, quantised to
\texttt{int8} and loaded into a BPF map; the forward pass uses integer
arithmetic and a fixed-point sigmoid lookup table with bounded iteration.
Features are restricted to quantities cheaply computable in hook context:
syscall number, hashed path prefix, descriptor type, destination port,
transaction depth and a short call n-gram. The classifier is \emph{fail-closed}:
an unknown or low-confidence result is treated as irrevocable and escalated,
never silently permitted (\S\ref{sec:threat}).

% =====================================================================
\section{Policy}
\label{sec:policy}

A deferral mechanism is only as meaningful as the rule governing release. This
section addresses the questions the mechanism alone leaves open.

\subsection{What is a transaction?}
\label{sec:boundary}

Nothing in \S\ref{sec:design} says how long a transaction lasts, and the answer
determines whether the guarantee is worth anything.

\begin{center}
\small
\begin{tabular}{@{}lp{4.1cm}@{}}
\toprule
\textbf{Boundary} & \textbf{Consequence}\\
\midrule
Per tool call & Guarantee window of milliseconds. Effects commit almost
                immediately; protection is near-nil\\
Per verification & Window equal to an edit--build--test cycle. Bounded
                staleness, bounded log\\
Per task or session & Large window, but unbounded log across an unbounded
                number of agent decisions, and a snapshot arbitrarily stale\\
\bottomrule
\end{tabular}
\end{center}

The question is not new but has never been answered. TxOS raised it and
deliberately left it open, observing that ``allowing a long running transaction
to isolate all system resources indefinitely is undesirable, so the OS may want
a policy that limits the size of a transaction,'' and comparing the matter to
controlling any process that abuses system resources
(\cite[\S3.1.3]{porter2009txos}). It supplied the mechanism and declined the
policy. For a workload whose transactions are bounded by human-scale activity
rather than by a critical section, the policy can no longer be deferred.

We argue for the middle option on principle rather than convenience. The reason
to speculate at all is that safety cannot be decided in advance; the moment
evidence \emph{does} arrive --- the build succeeds, the test suite passes, a
human reviews the diff --- the decision should be taken. A transaction should
therefore span exactly the interval over which no evidence is available. We call
this a \emph{verification-delimited} transaction. Its length is determined by
the workload rather than configured, and a session becomes a sequence of
transactions rather than one long one.

Two properties follow. First, the taxonomy and the boundary interlock: an
irrevocable effect terminates the enclosing transaction, exactly as an
irrevocable operation forces commitment in
STM~\cite{welc2008irrevocable}. Second, because agent transactions are orders of
magnitude longer than those the STM literature assumes, its results on
serialising irrevocable transactions cannot be imported unmodified; we expect
contention behaviour to differ and will report it.

We do not assert that verification-delimited is optimal. We make transaction
length an evaluation axis (\S\ref{sec:eval}) and report guarantee window,
staleness-induced conflict rate and log growth at each of the three boundaries.

\subsection{Who may commit}
\label{sec:authority}

The agent is the untrusted party. It must therefore not hold commit authority:
\texttt{tx\_commit} invoked by a process inside the transaction returns
\texttt{-EPERM}. Commit is exercised only by the supervising process, which is
outside the transaction and is driven by the verification signal, a policy rule,
or a human.

This single invariant eliminates the most obvious attack on the system --- a
prompt-injected agent instructed to commit before its exfiltration is
inspected --- and it is why \texttt{tx\_begin} and \texttt{tx\_commit} are
asymmetric in a way TxOS's interface was not.

\subsection{Coherence of agent reasoning}
\label{sec:coherence}

Reporting success for a deferred send makes the agent's subsequent reasoning
conditional on an event that may never occur. If the transaction later aborts,
the agent holds a chain of inferences built on a false premise.

TxOS faced the analogous problem and left it to the caller: on abort ``the OS
restores the kernel state to its pre-transaction state, but it does not revert
application state'' (\cite[\S3.1.4]{porter2009txos}), offering an optional
copy-on-write checkpoint of the application address space, in the manner of
Speculator, as a flag to \texttt{sys\_xbegin}. For an agent this cannot be
optional, and the address space is the wrong granularity in any case: the state
that must be reverted is the conversation, which lives in the supervising
harness rather than in the transacting process.

We therefore place the agent's reasoning state \emph{inside} the transaction.
Conversation and scratch state are owned by the supervising harness, staged in
the same copy-on-write layer as the filesystem, and discarded on abort along
with everything else. An abort therefore returns the agent to the
epistemic position it occupied at \texttt{tx\_begin}, with no false premise
surviving. This mirrors the versioned and transactional treatment of agent
memory developed in~\cite{chronomem2026,memtx2026}, and it is the reason our
transaction spans three kinds of state rather than two.

We considered the alternative of returning an explicit \texttt{DEFERRED} status
so the agent may reason about pending effects. It requires agent cooperation and
a model that handles the concept correctly, and we regard it as future work.

\subsection{Threat model}
\label{sec:threat}

The adversary is a prompt injection that fully controls the agent's outputs. It
does not hold a kernel exploit, does not control the supervising harness, and
does not run as root. Attacks on \sys{} itself and our responses:

\begin{center}
\small
\begin{tabular}{@{}p{2.6cm}p{4.6cm}@{}}
\toprule
\textbf{Attack} & \textbf{Response}\\
\midrule
Premature commit & Agent cannot commit (\S\ref{sec:authority})\\
Escape the transaction by spawning outside the namespace &
  Namespace is inherited across \texttt{fork}/\texttt{execve};
  \texttt{unshare} and \texttt{setns} denied inside a transaction\\
Adversarial evasion of the classifier &
  Fail-closed: low confidence escalates rather than permits. Evasion yields a
  prompt, not a silent pass\\
WAL exhaustion (emit millions of effects to force a flush) &
  Bounded log; overflow aborts the transaction rather than releasing it\\
TOCTTOU between classification and emission &
  Classification and disposition occur in the same hook invocation\\
Bypass TLS uprobes via static linking or custom crypto &
  Degrades to destination-only classification; an acknowledged coverage gap\\
Escape via un-instrumented paths &
  The general weakness of interception designs~\cite{sandboxescapes};
  we enumerate covered hooks explicitly rather than claim completeness\\
\bottomrule
\end{tabular}
\end{center}

Out of scope: kernel exploits, a compromised supervisor, malicious agent
binaries as opposed to injected agent \emph{instructions}, and timing side
channels --- deferral itself perturbs observable timing, which we note but do
not defend against.

% =====================================================================
\section{Implementation Plan}

The system decomposes into four components with interfaces frozen in week two
and one developer each. Every component ships a stub implementation on day one,
so any developer can build and boot the complete system in isolation and no
component is blocked on another's progress.

\begin{center}
\small
\begin{tabular}{@{}llp{2.4cm}@{}}
\toprule
& \textbf{Component} & \textbf{Kernel surface}\\
\midrule
P1 & Transaction core & syscall table,
                        \texttt{task\_struct}, locking\\
P2 & CoW storage & VFS, overlayfs, stackable FS\\
P3 & Effect interception & BPF LSM, networking, ring buffers\\
P4 & Policy, evaluation & verifier, quantisation, measurement\\
\bottomrule
\end{tabular}
\end{center}

Development proceeds in four independently demonstrable rungs: a loadable module
with \texttt{ioctl} control and overlay-backed abort; LSM interception with a
static rule table; replacement of that table by the in-kernel classifier; and
promotion to real system calls. Should a rung prove intractable, the preceding
rung remains a complete system. Multi-agent optimistic concurrency control is
scoped as a stretch objective.

% =====================================================================
\section{Evaluation Plan}
\label{sec:eval}

\paragraph{Project gate.} In weeks three to four we instrument real agent
sessions and measure the fraction of outbound network operations that are
fire-and-forget rather than request--response. Above 20\%, Contribution~1 stands
and the measurement becomes a headline figure. Below 10\%, we narrow the claim
and promote Contributions~2 and~5. This costs four days and de-risks the
remainder of the project.

\paragraph{Boundary sweep.} We run the full workload at each of the three
transaction boundaries of \S\ref{sec:boundary}, reporting guarantee window,
staleness-induced conflict rate at commit, peak WAL size, and the fraction of
transactions terminated early by an irrevocable effect. This converts the
central design decision from an assertion into a measurement.

\paragraph{Metrics.} Task success rate; containment rate under adversarial
input; interception latency; end-to-end overhead; storage amplification;
classifier accuracy before and after quantisation, and under adversarial
perturbation; and the empirical distribution of agent effects across the
taxonomy.

\paragraph{Baselines.} Unprotected execution; session checkpoints as shipped in
production agents~\cite{ccrewind}; \texttt{git} checkpoint and reset; container
rebuild; microVM snapshot; static \texttt{seccomp} allowlist; user-space policy
daemon via \texttt{seccomp} user-notify; and, where artifacts permit,
\cite{ftsandbox2025} and \cite{deltabox2026}, whose published overheads of
14.5\% and under 3\% are the figures we must meet or explain.

\paragraph{Workloads.} A benign corpus of realistic development tasks; a
fault-injection corpus of destructive actions; and an adversarial corpus in
which prompt injections are planted in files the agent reads during ordinary
operation, including the exfiltration scenario of \S1.

\paragraph{Headline figures.} (1) The deferrable fraction of agent outbound
effects. (2) Guarantee window against staleness across transaction boundaries.
(3) In-kernel classification latency against accuracy.

% =====================================================================
\section{Limitations}

Terminal output cannot be staged. Deferring a send means reporting success for
an operation that has not occurred; \S\ref{sec:coherence} addresses the
consequences for agent reasoning but not for a program that inspects the network
result directly. Coverage is the weakest point of any interception design.
Rollback mechanisms are themselves an attack surface~\cite{acrfence2026}. Our
classifier is trained on traces from a specific set of agents and may not
transfer. The compensable class depends on a registry somebody must maintain.
And request--response effects, which may prove to be the majority, are outside
the mechanism entirely.

% =====================================================================
\section{Conclusion}

System transactions were proposed in 2009 and set aside for want of a workload
justifying their cost. Autonomous agents supply that workload. Filesystem
rollback for agents is now solved and we do not contest it; what remains is the
class of effects the field has conceded as irreversible. We argue that class has
been drawn too broadly, because irreversibility is a property of observation
rather than of the system call, and that the kernel occupies the one interval in
which a guarantee is still available. Equally, we argue that the mechanism is
inert without an answer to what a transaction \emph{is} for an agent, and we
propose to answer that question by measurement rather than assertion.

\bibliographystyle{plain}
\bibliography{refs}

\end{document}
